% =============================================================================
%  CHAPITRE 4 : ANALYSE, DIAGNOSTIC ET INTERVENTION PROPOSÉE
% =============================================================================
\entete{Chapitre 4 : Diagnostic et intervention proposée}
\chapter{La restitution numérique de la visite à la Fondation Jean Félicien
Gacha : diagnostic et intervention}
\label{chap:diagnostic}

\introchap

Les données et les mesures rassemblées au chapitre précédent ne parlent pas
d'elles-mêmes : elles demandent à être interprétées. Ce chapitre procède donc au
diagnostic de la situation étudiée. Il analyse successivement chacune des quatre
variables explicatives (section~\ref{sec:analyse}), vérifie les hypothèses
formulées en introduction (section~\ref{sec:verif}), et dégage le modèle
explicatif qui se dégage des résultats (section~\ref{sec:modele}). Fort de ce
diagnostic, il présente ensuite l'intervention proposée, ses objectifs, ses
composantes, sa planification et sa faisabilité (section~\ref{sec:intervention}).

% =============================================================================
\section{Présentation et analyse de la situation}
\label{sec:analyse}

\subsection{Variable 1. La restitution spatiale : ce que l'espace exige et ce qu'il pardonne}

Les mesures relatives à la restitution spatiale conduisent à un constat en deux
temps.

D'une part, \textbf{la sensation d'espace ne dépend pas de la richesse du
modèle, mais de la continuité du geste}. Le moteur de panorama écrit pour ce
projet ne représente aucune géométrie : il calcule, pour chaque pixel, la
direction du rayon correspondant. Ce dispositif, techniquement modeste, produit
une perception d'espace supérieure à celle d'une galerie de photographies, parce
qu'il conserve la continuité du regard. À l'inverse, les imperfections de
raccord, la couture visible au bord du panorama, les facettes d'une sphère
grossièrement maillée, détruisent immédiatement l'illusion. L'espace pardonne
la simplification ; il ne pardonne pas la discontinuité.

D'autre part, \textbf{la perception des dimensions est un phénomène distinct de
la perception de l'espace}, et elle relève d'un autre dispositif. Aucune
rotation d'un modèle à l'écran ne renseigne sur la taille véritable d'un objet :
l'écran n'a pas d'échelle. Seule la projection dans l'environnement réel du
visiteur la restitue. C'est pourquoi la réalité augmentée n'a pas été traitée
comme un supplément d'agrément mais comme la seule réponse possible à une
question précise : \emph{quelle taille cela fait-il ?} Le refus de
l'agrandissement libre, au profit d'une échelle fixe mesurée sur le maillage,
découle directement de cette analyse.

Ce constat éclaire enfin une difficulté rencontrée : la couverture des appareils
n'est pas une question de programmation mais de \emph{format}. Un appareil Apple
ne déclenchera jamais de session de réalité augmentée à partir d'un fichier
GLB, quelle que soit la qualité du code. La couverture s'obtient en produisant
les deux formats, non en écrivant davantage de logique.

\subsection{Variable 2. La mise en relation documentaire : le vocabulaire comme obstacle}

C'est ici que se situe le résultat le plus instructif de l'étude. La mesure
comparative rapportée au chapitre~3 est reprise au tableau~\ref{tab:comparrech}.

\begin{table}[htbp]
\centering
\caption{Effet de la reformulation de la requête sur la recherche documentaire}
\label{tab:comparrech}
\small
\begin{tabularx}{\textwidth}{|X|C{3.0cm}|C{3.4cm}|}
\hline
\rowcolor{keycetete}
\thead{Formulation de la requête} & \thead{Résultats retournés} & \thead{Meilleur score obtenu} \\ \hline
Saisie française brute & 16 & 57 / 100 \\ \hline
Termes reformulés dans le vocabulaire de catalogage & 05 & \textbf{85 / 100} \\ \hline
\end{tabularx}
\source{Mesure conduite sur les mêmes collections et au même seuil de notation,
août 2026.}
\end{table}

L'interprétation de ce résultat mérite d'être exposée avec soin, car elle est
contre-intuitive. \textbf{Le nombre de résultats a chuté de plus des deux tiers,
et c'est une amélioration.} La recherche française ramenait un volume important
de correspondances faibles, tandis que la recherche reformulée fait
remonter, au sommet du classement, un masque éléphant conservé à Cleveland et
deux masques comparables du Metropolitan Museum, tous identifiés par leur numéro
d'inventaire. L'enrichissement échange donc du \emph{rappel} contre de la
\emph{précision}.

La portée de ce constat dépasse le cadre technique. Il établit que
\textbf{l'échec initial d'une recherche documentaire ne traduisait pas une
absence de données, mais un décalage de vocabulaire}. Les objets africains
dispersés dans les collections du monde sont bel et bien décrits, ils le sont
dans la langue et selon les conventions de catalogage des institutions qui les
conservent. Une institution camerounaise qui interroge ces catalogues dans ses
propres termes ne trouve rien, et pourrait en conclure à tort que ses œuvres
n'ont pas d'équivalent connu. Le décalage de vocabulaire fonctionne ainsi comme
une seconde dispersion, documentaire celle-là, qui se superpose à la dispersion
physique des objets.

La mesure conduite sur la langue des plongements confirme ce diagnostic sous un
autre angle : en anglais, les correspondances légitimes se séparent nettement du
bruit (au moins 0,898 contre au plus 0,779, soit une étendue de 0,20), tandis
qu'en français cette séparation se réduit à 0,13, au point qu'une peinture sans
rapport pouvait être classée devant deux masques réellement apparentés. Ce n'est
pas la méthode qui échouait, c'est la langue dans laquelle on la sollicitait.

Une limite doit toutefois être rappelée : la comparaison porte sur le
\emph{texte} des notices et non sur l'\emph{image} des œuvres. Deux objets
visuellement très proches mais décrits selon des conventions différentes peuvent
donc échapper au rapprochement. Cette limite est structurelle, non accidentelle,
et figure au premier rang des perspectives.

\subsection{Variable 3. L'ancrage documentaire : ce que l'assistant refuse de dire}

L'analyse des résultats relatifs à l'assistant conduit à déplacer le critère
d'évaluation. On juge spontanément un assistant à ce qu'il sait ; l'étude montre
qu'en médiation culturelle il faut le juger à \textbf{ce qu'il refuse
d'affirmer}.

Trois mécanismes se sont révélés déterminants, et leur ordre d'efficacité est
inverse de leur complexité technique.

\begin{enumerate}
  \item \textbf{Le filtrage en amont}, le plus simple, est le plus efficace.
        Écarter les questions manifestement hors périmètre avant même de
        solliciter le modèle de langue supprime une classe entière d'erreurs, à
        coût nul.

  \item \textbf{La restriction du périmètre documentaire} au musée ou à la salle
        où se trouve le visiteur réduit mécaniquement le risque de confusion
        entre deux œuvres homonymes ou similaires.

  \item \textbf{L'autorisation explicite d'ignorer}, enfin, est le mécanisme le
        moins intuitif et le plus important. Un modèle de langue à qui l'on
        demande une réponse en produira une. Il faut donc lui accorder
        formellement le droit de n'en pas produire. Sans cette clause, la
        consigne d'ancrage documentaire reste insuffisante : le modèle comble
        les vides.
\end{enumerate}

Le journal des questions restées sans réponse apporte un enseignement
complémentaire, d'ordre organisationnel. Il transforme une défaillance de
l'assistant en indication utile pour l'institution : une question récurrente
sans réponse ne signale pas un défaut du dispositif, elle signale une lacune de
la documentation. Le système cesse alors d'être seulement un outil de diffusion
pour devenir un outil de pilotage documentaire.

\subsection{Variable 4. Le cloisonnement : pourquoi la garantie doit être portée par la base}

Les résultats relatifs à la quatrième variable établissent une distinction que
l'expérience du projet a rendue tangible : celle entre une garantie
\emph{applicative} et une garantie \emph{structurelle}.

Une garantie applicative consiste à ce que le programme n'aille pas chercher les
données d'autrui. Elle repose sur la correction de chaque requête écrite, et
elle tombe dès qu'une seule d'entre elles est fautive. Deux incidents survenus
au cours du projet l'ont confirmé : un classement des œuvres les plus consultées
qui ne filtrait pas sur l'institution affichée, et un droit d'accès valable
au-delà de l'institution émettrice. Dans les deux cas, le code demandait
sincèrement ce qu'il croyait devoir demander, et la fuite s'est produite.

Une garantie structurelle consiste à ce que la base de données refuse de
répondre. Elle ne dépend pas de la correction du programme. C'est celle qui a
été retenue : chaque table métier porte une règle de sécurité au niveau de la
ligne, et le rattachement à l'institution est posé automatiquement à
l'insertion. Le même principe a été appliqué à l'assistant d'installation, qui
écrit avec le jeton de la personne qui l'utilise et non avec une clé
privilégiée : un agent détourné par une consigne malveillante se heurte au refus
de la base, et non à la vigilance du prompt.

Ce diagnostic a une portée qui excède le projet : \textbf{dans une architecture
accueillant plusieurs organisations, l'isolement ne peut pas être une propriété
du code ; il doit être une propriété de la donnée.}

\subsection{Analyse d'ensemble : forces et faiblesses}

\begin{table}[htbp]
\centering
\caption{Analyse des forces, faiblesses, opportunités et menaces du dispositif}
\label{tab:swot}
\footnotesize
\begin{tabularx}{\textwidth}{|X|X|}
\hline
\rowcolor{keycetete}
\thead{Forces} & \thead{Faiblesses} \\ \hline
\begin{minipage}[t]{\linewidth}\vspace{2pt}
\begin{itemize}[leftmargin=1.1em,topsep=0pt]
  \item Ensemble cohérent réunissant gestion, immersion et assistance
  \item Cloisonnement garanti par la base de données
  \item Coût récurrent d'hébergement très faible (730 FCFA/mois)
  \item Aucune dépendance externe pour les moteurs de rendu
  \item Démonstration possible sur une base vierge
  \item Chaîne de déploiement automatisée et vérifiée
\end{itemize}\vspace{2pt}\end{minipage}
&
\begin{minipage}[t]{\linewidth}\vspace{2pt}
\begin{itemize}[leftmargin=1.1em,topsep=0pt]
  \item Comparaison portant sur le texte et non sur l'image
  \item Quatre grandes collections inaccessibles faute de clé
  \item Aucune carte de profondeur produite en production
  \item Essais de réalité augmentée sur peu d'appareils
  \item Dépendance à des paliers gratuits de services tiers
  \item Contenu numérique de la Fondation encore à produire
\end{itemize}\vspace{2pt}\end{minipage} \\ \hline
\rowcolor{keycetete}
\thead{Opportunités} & \thead{Menaces} \\ \hline
\begin{minipage}[t]{\linewidth}\vspace{2pt}
\begin{itemize}[leftmargin=1.1em,topsep=0pt]
  \item Tissu dense d'institutions patrimoniales dans l'Ouest
  \item Coût par institution décroissant grâce au catalogue mondial partagé
  \item Intérêt de la diaspora pour le patrimoine d'origine
  \item Ouverture progressive des collections mondiales
  \item Généralisation des capteurs LiDAR sur téléphone
\end{itemize}\vspace{2pt}\end{minipage}
&
\begin{minipage}[t]{\linewidth}\vspace{2pt}
\begin{itemize}[leftmargin=1.1em,topsep=0pt]
  \item Mise en sommeil des services gratuits faute de trafic
  \item Évolution ou fermeture des interfaces des collections
  \item Départ de la personne détentrice du savoir technique
  \item Connectivité limitée pour une partie du public visé
  \item Coût de production des médias (prises de vue, numérisation)
\end{itemize}\vspace{2pt}\end{minipage} \\ \hline
\end{tabularx}
\source{Élaboré par nos soins à partir du diagnostic.}
\end{table}

% =============================================================================
\section{Vérification des hypothèses}
\label{sec:verif}

Le tableau~\ref{tab:hypotheses} confronte chaque hypothèse aux résultats
obtenus.

\begin{table}[htbp]
\centering
\caption{Vérification des hypothèses de l'étude}
\label{tab:hypotheses}
\footnotesize
\begin{tabularx}{\textwidth}{|C{2.1cm}|X|C{2.6cm}|}
\hline
\rowcolor{keycetete}
\thead{Hypothèse} & \thead{Éléments de vérification} & \thead{Conclusion} \\ \hline
\textbf{Spécifique 1} &
Le parcours à 360 degrés restitue le déplacement de salle en salle ; la
présentation tridimensionnelle permet l'examen sous tous les angles ; la réalité
augmentée restitue les dimensions véritables (45 $\times$ 52 $\times$ 45 cm
mesurés sur le maillage) sur les trois familles d'appareils visées. La
perception des dimensions, en revanche, n'est obtenue que par la réalité
augmentée, et non par la rotation à l'écran. &
\textbf{Confirmée}, avec la précision que deux dispositifs distincts sont
requis \\ \hline
\textbf{Spécifique 2} &
Cinq collections ouvertes interrogées ; 41 œuvres apparentées identifiées ;
meilleur score de correspondance de 90 sur 100 ; 137 pays représentés dans la
dispersion. Le gain de précision apporté par la reformulation est mesuré
(57 $\rightarrow$ 85). &
\textbf{Confirmée} \\ \hline
\textbf{Spécifique 3} &
Les trois garde-fous, filtrage en amont, restriction du périmètre, autorisation
explicite d'ignorer, produisent un assistant qui refuse les sujets hors
périmètre, cite ses sources et déclare son ignorance. Le journal des questions
sans réponse objective les limites restantes. &
\textbf{Confirmée} \\ \hline
\textbf{Spécifique 4} &
38 tables protégées par des règles de sécurité au niveau de la ligne ;
rattachement automatique à l'insertion ; accueil d'une nouvelle institution sans
modification du code ; deux fuites détectées et corrigées, l'une et l'autre
d'origine applicative et non structurelle. &
\textbf{Confirmée}, la garantie devant être portée par la base \\ \hline
\textbf{Générale} &
Les quatre hypothèses spécifiques étant vérifiées, et le dispositif étant en
service, l'hypothèse générale est corroborée. La restitution demeure toutefois
tributaire de la production des médias par l'institution. &
\textbf{Confirmée sous réserve} \\ \hline
\end{tabularx}
\source{Élaboré par nos soins à partir des mesures du chapitre~3.}
\end{table}

% =============================================================================
\section{Modèle explicatif issu des résultats}
\label{sec:modele}

L'ensemble du diagnostic peut être ramené à une proposition unique :
\textbf{la restitution numérique d'une visite repose sur quatre continuités,
dont la rupture de l'une suffit à faire échouer l'ensemble}.

\begin{itemize}
  \item \textbf{La continuité spatiale} : le visiteur doit pouvoir passer d'un
        point de vue au suivant sans rupture. Une galerie d'images la brise.
  \item \textbf{La continuité d'échelle} : l'objet doit conserver ses
        dimensions. L'écran, dépourvu d'échelle, la brise ; la réalité augmentée
        la rétablit.
  \item \textbf{La continuité documentaire} : chaque affirmation doit être
        rattachable à une source. L'invention la brise ; l'ancrage documentaire
        la maintient.
  \item \textbf{La continuité narrative} : chaque écran doit permettre de
        poursuivre le fil, de l'objet au chef, du chef à la lignée, de la
        lignée aux œuvres apparentées. Un écran sans issue la brise.
\end{itemize}

% =============================================================================
\section{Intervention proposée et justification}
\label{sec:intervention}

\subsection{Présentation de l'intervention}

L'intervention proposée consiste à \textbf{faire passer la plateforme MUSÉA de
l'état de démonstration éprouvée à celui de service exploité durablement par la
Fondation}. Elle comporte cinq composantes : l'industrialisation de la chaîne de
production par la conteneurisation, la consolidation de la chaîne documentaire
de l'assistant, la campagne de production des médias, le transfert de
compétences, et l'ouverture aux institutions partenaires.

\subsection{Justification de l'intervention}

Trois constats issus du diagnostic la rendent nécessaire.

\begin{enumerate}
  \item \textbf{L'architecture actuelle atteint sa limite fonctionnelle.} Le
        dispositif d'hébergement retenu, un stockage d'objets servi par un
        réseau de diffusion, convient parfaitement à un site statique et son
        coût est dérisoire. Mais plusieurs traitements identifiés comme
        nécessaires ne peuvent s'y exécuter : l'ingestion périodique des
        documents dans l'index vectoriel, la reconstruction photogrammétrique,
        la production des cartes de profondeur, et l'interrogation programmée
        des collections. Ces traitements exigent un environnement d'exécution
        durable, que seul un service de conteneurs peut fournir à coût maîtrisé.

  \item \textbf{La reprise par la Fondation suppose la reproductibilité.} Tant
        que l'environnement de production diffère de l'environnement de
        développement, la reprise reste risquée. Le conteneur supprime cet écart
        en transportant l'environnement complet et non le seul code.

  \item \textbf{La pérennité suppose la maîtrise du coût.} Passer à des
        conteneurs augmente mécaniquement la dépense. Cette augmentation doit
        donc être instruite, plafonnée et suivie, d'où la place accordée au
        FinOps dans la présente proposition.
\end{enumerate}

\subsection{Objectifs d'intervention}

\encadre{Objectif général d'intervention}{%
Mettre la plateforme MUSÉA en exploitation durable au sein de la Fondation Jean
Félicien Gacha, en industrialisant sa chaîne de production, en complétant son
contenu numérique et en transférant sa maîtrise au personnel de l'institution.}

\begin{itemize}[label={},leftmargin=0pt,itemsep=5pt,topsep=4pt,parsep=0pt]
  \item \textbf{Objectif d'intervention 1.} Industrialiser la chaîne de production par la
        conteneurisation et l'orchestration, avec des environnements de
        pré-production et de production distincts.
  \item \textbf{Objectif d'intervention 2.} Consolider la chaîne documentaire de l'assistant, de
        l'ingestion des documents à l'indexation vectorielle.
  \item \textbf{Objectif d'intervention 3.} Produire les médias manquants : prises de vue à
        360~degrés de l'ensemble des espaces et numérisation d'un premier
        ensemble d'œuvres.
  \item \textbf{Objectif d'intervention 4.} Transférer au personnel la maîtrise de l'outil au moyen
        d'une formation et d'une documentation d'exploitation.
  \item \textbf{Objectif d'intervention 5.} Accueillir deux institutions partenaires afin d'éprouver
        l'ouverture de l'architecture en conditions réelles.
\end{itemize}

\subsection{Composante 1 : Industrialisation par la conteneurisation}

\subsubsection{Principe de la chaîne cible}

La chaîne proposée conserve le principe déjà en service, toute modification
déposée est vérifiée, puis publiée automatiquement, mais substitue à la
publication de fichiers la construction et le déploiement d'une \emph{image de
conteneur}. Son enchaînement est le suivant :

\begin{enumerate}
  \item le développeur dépose sa modification sur le dépôt de code ;
  \item l'intégration continue installe les dépendances, compile l'application
        et exécute les contrôles de non-régression ;
  \item une \textbf{image Docker} multi-étapes est construite : une première
        étape compile l'application, une seconde ne conserve que le résultat et
        le serveur de diffusion, ce qui produit une image légère et dépourvue
        d'outils de compilation ;
  \item l'image est signée et déposée dans le \textbf{registre d'images} ;
  \item le déploiement en \textbf{pré-production} est automatique ; les essais
        de bout en bout y sont exécutés ;
  \item le déploiement en \textbf{production} intervient après validation, par
        remplacement progressif des instances, sans interruption de service ;
  \item la vérification a posteriori du service en ligne conditionne la
        clôture ; en cas d'échec, le retour à la version précédente est
        automatique.
\end{enumerate}

\begin{figure}[htbp]
\centering
\schema{devops-cible}
\caption{Chaîne DevOps cible : de la modification du code à la production par conteneurs}
\label{fig:devops-cible}
\source{Élaboré par nos soins.}
\end{figure}

\subsubsection{Choix de l'orchestrateur}

Le choix entre les deux orchestrateurs disponibles a été instruit ; ses termes
sont présentés au tableau~\ref{tab:orchestrateur}, reporté en annexe~11. La
comparaison porte sur le coût fixe, la charge d'administration, la courbe
d'apprentissage et la capacité de montée en charge.

\encadre{Recommandation}{%
\textbf{Retenir ECS avec Fargate dans un premier temps.} Pour une institution
disposant d'une équipe technique réduite, le coût fixe et la charge
d'administration de Kubernetes ne se justifieraient pas au volume actuel. La
description de l'infrastructure par le code et la conteneurisation préservent
toutefois intégralement la possibilité d'une migration ultérieure vers EKS :
c'est la même image qui serait déployée. La bascule devra être envisagée le
jour où la plateforme dépassera une dizaine d'institutions ou exigera plusieurs
services distincts fonctionnant en permanence.}

\subsection{Composante 2 : Consolidation de la chaîne documentaire}

Trois actions sont proposées.

\begin{enumerate}
  \item \textbf{Automatiser l'ingestion.} L'export des notices vers le dépôt
        documentaire, leur découpage, leur normalisation en anglais et leur
        vectorisation doivent devenir une tâche planifiée, exécutée
        quotidiennement dans un conteneur. Aujourd'hui, cette chaîne est
        déclenchée manuellement.

  \item \textbf{Éprouver la chaîne complète de rapprochement en conditions
        réelles.} Chaque maillon a été vérifié séparément ; l'enchaînement
        complet (mise en cache, vectorisation, reclassement, proposition) n'a jamais été exécuté de bout en bout avec un compte du personnel. Cet
        essai constitue une priorité immédiate.

  \item \textbf{Obtenir les clés d'accès manquantes.} Les demandes auprès
        d'Europeana, du Rijksmuseum, de la Smithsonian Institution et des
        Harvard Art Museums doivent être introduites : ces quatre sources
        élargiraient sensiblement la couverture, la liste blanche technique
        étant déjà prête à les accueillir.
\end{enumerate}

\subsection{Composante 3 : Production des médias}

Le dispositif est prêt ; le contenu reste à produire. Cette composante prévoit
une campagne de prises de vue à 360~degrés couvrant l'ensemble des espaces, la
numérisation d'un premier ensemble d'œuvres emblématiques par capteur LiDAR et
par photogrammétrie, la production des versions au format USDZ nécessaires aux
appareils Apple, ainsi que la génération des cartes de profondeur pour les
œuvres non numérisées en trois dimensions. Elle prévoit également
l'enregistrement ou la synthèse des récits associés aux pièces principales.

\subsection{Composante 4 : Transfert de compétences}

Un dispositif qu'une seule personne sait exploiter n'est pas un dispositif
pérenne. Cette composante prévoit une formation du personnel à l'administration
courante, une formation à la production des médias, la remise d'une
documentation d'exploitation, et la désignation d'un référent interne.

\subsection{Composante 5 : Ouverture aux institutions partenaires}

L'ouverture de l'architecture ne sera véritablement établie qu'après avoir été
éprouvée. Cette composante prévoit l'accueil de deux institutions partenaires,
la mesure du délai réel d'installation, et le recueil de leurs retours.

\subsection{Planification de l'intervention}

L'intervention se déroule sur douze mois, selon le calendrier prévisionnel
présenté au tableau~\ref{tab:planning}, reporté en annexe~10. Les deux
premiers trimestres portent l'industrialisation et la consolidation
documentaire, qui conditionnent les composantes suivantes ; la production des
médias et le transfert de compétences se poursuivent ensuite en continu, et
l'ouverture aux institutions partenaires n'intervient qu'une fois la chaîne
éprouvée.

\subsection{Étude de faisabilité}

\paragraph{Sur le plan technique, la faisabilité est acquise.} L'ensemble des briques nécessaires existe et a été éprouvé :
l'image de conteneur est déjà décrite et sert d'ores et déjà à reproduire en
local le comportement de production ; l'authentification sans clé permanente est
en service ; l'infrastructure est décrite par le code. Le passage aux conteneurs
consiste à substituer une cible de déploiement à une autre, non à reconstruire
la chaîne.

\paragraph{Sur le plan économique,} le tableau~\ref{tab:coutfonctionnement}
présente le coût de fonctionnement annuel projeté dans l'hypothèse de la chaîne
cible.

\begin{table}[htbp]
\centering
\caption{Coût de fonctionnement annuel projeté après intervention}
\label{tab:coutfonctionnement}
\small
\begin{tabularx}{\textwidth}{|L{4.6cm}|X|R{2.6cm}|}
\hline
\rowcolor{keycetete}
\thead{Poste} & \thead{Hypothèse retenue} & \thead{Montant annuel (FCFA)} \\ \hline
Hébergement statique et diffusion & Stockage, réseau de diffusion, noms de domaine, usage constaté & 8 760 \\ \hline
Nom de domaine & Renouvellement annuel & 9 000 \\ \hline
Exécution en conteneurs & Deux services de taille minimale, facturation à la ressource & 96 000 \\ \hline
Registre d'images & Stockage des images et analyse de vulnérabilités & 6 000 \\ \hline
Journalisation et supervision & Conservation des journaux sur 30 jours & 12 000 \\ \hline
Calcul graphique ponctuel & Photogrammétrie, 40 heures par an & 12 600 \\ \hline
Base de données et fonctions serveur & Palier payant recommandé pour éviter la mise en sommeil & 165 000 \\ \hline
Génération de texte et synthèse vocale & Paliers d'entrée & 60 000 \\ \hline
\rowcolor{keycetotal}
\textbf{Total de fonctionnement annuel} & & \textbf{369 360} \\ \hline
\end{tabularx}
\source{Projection établie par nos soins à partir des tarifs publics des
fournisseurs et de la consommation constatée durant le stage.}
\end{table}

Ce montant appelle deux observations. D'une part, il reste inférieur au coût
d'un poste de travail unique, ce qui le place dans l'ordre de grandeur
soutenable pour une institution de la taille de la Fondation. D'autre part, il
est \textbf{mutualisable} : le catalogue mondial des objets externes étant
partagé, l'arrivée d'institutions partenaires n'en augmente pas la partie la
plus coûteuse, mais en répartit la charge.

\paragraph{Sur le plan organisationnel}, la faisabilité constitue le point le plus délicat : elle suppose la désignation d'un
référent interne, la formation effective du personnel, et l'inscription de la
production des médias dans les activités courantes de la Fondation. La composante~4 y pourvoit, mais son succès dépend d'une décision institutionnelle
qui échappe au périmètre technique.

\subsection{Risques et mesures d'atténuation}

Sept risques ont été identifiés et assortis d'une mesure d'atténuation ; le
tableau~\ref{tab:risques}, reporté en annexe~12, les recense avec leur
probabilité et leur gravité. Deux d'entre eux commandent les autres. La mise en
sommeil de la base de données faute de trafic, probable et grave, se traite par
la souscription du palier payant ou, à défaut, par une sollicitation périodique
programmée. Le départ du référent technique, moins probable mais tout aussi
grave, appelle une documentation d'exploitation et la formation d'un second
référent : c'est le risque que la démarche DevOps décrite au chapitre~1 réduit
le plus directement, une mise en production décrite dans un fichier étant
transmissible là où une suite de gestes mémorisés ne l'est pas.

\conclchap

Ce chapitre a interprété les résultats de l'étude et vérifié les hypothèses
formulées en introduction : les quatre hypothèses spécifiques sont confirmées,
et l'hypothèse générale l'est sous la réserve que l'institution produise le
contenu numérique dont le dispositif a besoin. Le diagnostic a fait apparaître
deux enseignements qui dépassent le cas étudié : \emph{le décalage de
vocabulaire constitue une seconde dispersion des œuvres, documentaire celle-là},
et \emph{l'isolement des données ne peut être une propriété du code, il doit
être une propriété de la donnée}. Le modèle des quatre continuités en offre une
formulation synthétique. Sur cette base, une intervention en cinq composantes a
été proposée, planifiée, chiffrée et assortie de son analyse de risques. La
conclusion générale en tire les enseignements d'ensemble et ouvre les
perspectives.
